% ============================================================
% secciones/contenido.tex
% EXP-1 - Netflix: comunicación entre microservicios con RPC y gRPC,
% comparada con Java RMI
% Estructura de la exposición (12 min) según el enunciado.
% Cada sección: qué hay que cubrir, tiempo y quién la expone.
% El texto de cada sección lo escribe el integrante responsable.
% ============================================================

% ------------------------------------------------------------
\section{Presentación del sistema}
% Tiempo: 1 min | Expone: [NOMBRE]
% Cubrir:
%  - Qué problema resuelve Netflix (streaming de video a escala global).
%  - Quién lo usa y a qué escala: número de suscriptores y países
%    (citar el informe a inversionistas más reciente de Netflix).
%  - Por qué nos interesa: su comunicación interna entre microservicios
%    usa RPC, en particular gRPC.

% ------------------------------------------------------------
\section{Arquitectura}
% Tiempo: 2 min | Expone: [NOMBRE]
% Cubrir:
%  - Componentes: clientes (apps), gateway de entrada, microservicios
%    internos, almacenamiento.
%  - Cómo se comunican: llamadas RPC entre microservicios con gRPC
%    (HTTP/2 + Protocol Buffers).
%  - Evolución: antes usaban su propio stack sobre HTTP/1.1; luego
%    migraron a gRPC (fuente: caso de estudio de la CNCF sobre Netflix).
%  - Diagrama propio en draw.io (diagramas/) + esquema a mano (anexo).

% ------------------------------------------------------------
\section{Transparencias de distribución}
% Tiempo: 1,5 min | Expone: [NOMBRE]
% Cubrir (con un ejemplo concreto de cada una):
%  - Acceso: el stub de gRPC hace que llamar a un servicio remoto se vea
%    como llamar a una función local.
%  - Ubicación: el cliente no sabe en qué máquina está el servicio.
%  - Replicación y fallos: varias instancias de cada servicio.
%  - Cuáles NO ofrece y por qué (ej.: la latencia de red no se puede
%    ocultar del todo; un fallo de red sí puede verse como error/timeout).

% ------------------------------------------------------------
\section{CAP y PACELC}
% Tiempo: 2 min | Expone: [NOMBRE]
% Cubrir:
%  - Qué prioriza Netflix ante una partición: disponibilidad (AP).
%  - Qué prioriza sin partición: latencia (EL).
%  - Evidencia: documentación técnica o artículos (Netflix Tech Blog).
%  - Aclarar: gRPC es el mecanismo de comunicación; el CAP es del
%    sistema (Netflix), no del protocolo.

% ------------------------------------------------------------
\section{Modelos de comunicación}
% Tiempo: 1,5 min | Expone: [NOMBRE]
% Cubrir:
%  - Cliente-servidor entre microservicios.
%  - Sincrónica (llamada unaria) y streaming de gRPC.
%  - Mecanismo: RPC (gRPC) frente a sockets y mensajería.
%  - Por qué Netflix eligió gRPC: necesitaba clientes que funcionaran
%    entre lenguajes, sobre todo Java y Node.js (caso CNCF).
%  - Aquí entra Java RMI: solo funciona entre programas Java.

% ------------------------------------------------------------
\section{Tolerancia a fallos}
% Tiempo: 2 min | Expone: [NOMBRE]
% Cubrir:
%  - Qué fallos soporta: caída de instancias, de zonas.
%  - Cómo los detecta: timeouts/deadlines de gRPC, health checks.
%  - Cómo se recupera: reintentos, balanceo hacia otras instancias,
%    Chaos Monkey para probarlo.
%  - Incidente real documentado de Netflix + análisis de causa raíz
%    (método de FORM-AC-01: causa inmediata, causa raíz, mitigación).

% ------------------------------------------------------------
\section{Comparación y postura crítica}
% Tiempo: 1 min | Expone: [NOMBRE]
% Cubrir:
%  - Tabla: gRPC (Netflix) vs Java RMI, con al menos 5 criterios
%    (lenguajes, transporte, serialización, rendimiento, streaming,
%    tolerancia a fallos).
%  - Fortalezas, limitaciones y trade-offs de gRPC.
%  - Cuándo RMI todavía tendría sentido.

% ------------------------------------------------------------
\section{Lecciones para el PFC y conclusiones}
% Tiempo: 1 min | Expone: [NOMBRE]
% Cubrir:
%  - Qué decisión de Netflix aplicaríamos en nuestro PFC
%    (ej.: definir contratos entre servicios, timeouts y reintentos).
%  - 3 o 4 conclusiones concretas, no teoría repetida.
